%! Author = Eternal_Core
%! Date = 06.04.2026

\documentclass[12pt,a4paper]{article}

\usepackage[T2A]{fontenc}
\usepackage[utf8]{inputenc}
\usepackage[russian]{babel}
\usepackage{amsmath}
\usepackage{geometry}
\usepackage{listings}
\usepackage{xcolor}
\usepackage{hyperref}

\geometry{left=2.5cm,right=1.5cm,top=2cm,bottom=2cm}

\lstset{
    language=Python,
    basicstyle=\ttfamily\small,
    keywordstyle=\color{blue},
    stringstyle=\color{red},
    commentstyle=\color{gray},
    showstringspaces=false,
    breaklines=true,
    frame=single,
    tabsize=4,
    captionpos=b
}

\title{\textbf{Отчёт по лабораторной работе №7} \\ \large Работа с API}
\author{}
\date{}

\begin{document}

\maketitle
\thispagestyle{empty}

\section*{Цель работы}
Изучить принципы работы с REST API, научиться формировать HTTP-запросы с параметрами, получать и обрабатывать JSON-ответы, а также создать графический интерфейс для загрузки и отображения изображений из сети.

\section*{Задание 1: OpenWeatherMap API}
\subsection*{Описание}
Используя сервис OpenWeatherMap, реализовать программу, которая показывает температуру, ощущаемую температуру, влажность, давление, описание погоды и скорость ветра в указанном городе.

\subsection*{Реализация}
Для запросов используется библиотека \texttt{requests}. API-ключ передаётся в параметрах запроса. Базовый URL:
\begin{lstlisting}
https://api.openweathermap.org/data/2.5/weather
\end{lstlisting}

Параметры запроса:
\begin{itemize}
    \item \texttt{q} --- название города
    \item \texttt{appid} --- API-ключ
    \item \texttt{units=metric} --- температура в Цельсиях
    \item \texttt{lang=ru} --- описание на русском языке
\end{itemize}

\subsection*{Фрагмент кода}
\begin{lstlisting}
def get_weather(city_name):
    params = {
        "q": city_name,
        "appid": OWM_API_KEY,
        "units": "metric",
        "lang": "ru"
    }
    response = requests.get(OWM_BASE_URL, params=params, timeout=10)
    response.raise_for_status()
    data = response.json()
\end{lstlisting}

\subsection*{Результат}
Программа запрашивает название города, отправляет GET-запрос к API и выводит структурированную информацию: температура, ощущается как, влажность, давление, описание, скорость ветра. При неверном названии города функция возвращает \texttt{False} и выводит сообщение об ошибке.

\section*{Задание 2: Rick and Morty API}
\subsection*{Описание}
Использовать API сериала Rick and Morty для поиска персонажей и получения случайных персонажей. Вариант 9 --- \url{https://rickandmortyapi.com}.

\subsection*{Реализация}
API не требует авторизации. Для поиска персонажа используется эндпоинт:
\begin{lstlisting}
https://rickandmortyapi.com/api/character?name=<имя>
\end{lstlisting}

Для получения случайного персонажа сначала запрашивается общее количество персонажей через \texttt{/api/character}, затем генерируется случайный ID от 1 до \texttt{count} и запрашивается \texttt{/api/character/\{id\}}.

\subsection*{Фрагмент кода}
\begin{lstlisting}
def get_character(name):
    params = {"name": name}
    response = requests.get(BASE_URL, params=params, timeout=10)
    response.raise_for_status()
    data = response.json()
    char = data["results"][0]
\end{lstlisting}

\subsection*{Результат}
Программа выводит: имя, пол, вид, статус, родной мир и текущее место персонажа. При вводе пустого имени --- выбирается случайный персонаж.

\section*{Задание 3 (бонус): Генератор котиков}
\subsection*{Описание}
Используя \url{https://cataas.com} и библиотеку \texttt{tkinter}, создать приложение, отображающее случайную GIF-картинку с котом. Окошко состоит из картинки и кнопки, нажатие на которую загружает нового кота.

\subsection*{Реализация}
Сервис cataas.com предоставляет REST API для получения случайных изображений котов. Эндпоинт \texttt{/cat/gif} возвращает анимированную GIF-картинку. Приложение загружает бинарные данные изображения, извлекает отдельные кадры GIF с помощью \texttt{PIL.Image}, и анимирует их через \texttt{tkinter.after()} с учётом задержки каждого кадра.

\subsection*{Фрагмент кода}
\begin{lstlisting}
CATAAS_URL = "https://cataas.com/cat/gif"

def load_image(self):
    response = requests.get(CATAAS_URL, timeout=15)
    response.raise_for_status()
    gif = Image.open(BytesIO(response.content))
    if gif.format == "GIF":
        self._load_animated_gif(gif)
        self._play_animation()
\end{lstlisting}

\subsection*{Анимация GIF}
Библиотека tkinter не поддерживает анимацию GIF нативно. Для отображения анимации:
\begin{enumerate}
    \item GIF разделяется на отдельные кадры через \texttt{gif.seek()}
    \item Для каждого кадра извлекается задержка из \texttt{gif.info["duration"]}
    \item Кадры последовательно показываются через \texttt{root.after(delay, callback)}
\end{enumerate}

\section*{Интерфейс командной строки}
При запуске программы пользователю предлагается меню:
\begin{lstlisting}
Выберите задание:
  1 - Погода в городе
  2 - Персонаж Rick and Morty
  3 - Генератор котиков
  q - выход
\end{lstlisting}

После выполнения задания меню показывается снова. Выход --- клавиша \texttt{q}.

\section*{Используемые библиотеки}
\begin{itemize}
    \item \texttt{requests} --- HTTP-запросы к API
    \item \texttt{tkinter} --- графический интерфейс
    \item \texttt{PIL (Pillow)} --- обработка изображений, извлечение кадров GIF
    \item \texttt{io.BytesIO} --- работа с бинарными данными в памяти
    \item \texttt{random} --- генерация случайного ID персонажа
\end{itemize}

\section*{Вывод}
В ходе лабораторной работы были изучены принципы работы с REST API, полученные навыки формирования HTTP-запросов с параметрами, обработки JSON-ответов и ошибок сети, а также создано графическое приложение для загрузки и отображения анимированных изображений из сети.

\end{document}
